\section{Reliability and Evaluation Protocol}
\label{sec:eval}

\paragraph{Task metrics.}
We report macro-F1 for crop type, presence, count bin, density, and coverage,
plus exact schema validity.  A task macro-F1 averages the five task-level
macro-F1 values so that a dense count category cannot dominate the report.
All intervals and paired differences cluster by source image.  Tile-level
confidence intervals are not used because overlapping tiles are correlated.

\paragraph{Repeat consistency.}
For repeated generations $\{\hat y^{(r)}\}_{r=1}^{R}$ and categorical task set
$\mathcal{K}$, consistency is the mean modal agreement
\begin{equation}
C=\frac{1}{|\mathcal{K}|}\sum_{k\in\mathcal{K}}
\frac{1}{R}\max_{a}\sum_{r=1}^{R}\mathbf{1}[\hat y_k^{(r)}=a].
\label{eq:consistency}
\end{equation}
The final experiment uses three generations under identical decoding settings.

\paragraph{Risk score and abstention.}
The fixed diagnostic score combines schema validity $s$, repeat consistency
$C$, VLM confidence $p$, and detector agreement $d$:
\begin{equation}
u_{\mathrm{fixed}}=.35s+.25C+.20p+.20d.
\label{eq:fixed-gate}
\end{equation}
The proposed gate is a validation-fitted failure predictor using the same
signals.  Its threshold is selected on validation data and frozen before test
evaluation.  We report risk--coverage curves, area under the risk--coverage
curve (AURC), accepted risk, false-accept rate, and coverage.  Verified empty
sources define the hallucination denominator; ordinary detector misses do not
automatically become VLM hallucinations.

\paragraph{Cross-crop protocols.}
The complete matrix contains Magnolia$\rightarrow$Magnolia,
Azalea$\rightarrow$Azalea, Magnolia$\rightarrow$Azalea,
Azalea$\rightarrow$Magnolia, and pooled training/testing.  Model selection,
selector choice, risk threshold, and prompt are frozen on training/validation
sources for each protocol.  No test source contributes to threshold fitting.

\paragraph{Efficiency.}
Latency excludes download and reports model loading separately.  Runs record
warm-up, batch size, precision, generation limit, hardware, and peak incremental
memory.  Matched comparisons use the same budget and prompt; all-tile inference
is an efficiency baseline, not a matched-input accuracy baseline.
